\title{2023 年高考数学 全国甲卷-文科}
\date{}
\maketitle

\section*{一、单项选择题}
本题共12小题，每小题5分，共60分.在每小题给出的四个选项中，只有一项是符合题目要求的.

\begin{question}[points=5]
设全集 $U = \{1,2,3,4,5\}$，集合 $M = \{1,4\}, N = \{2,5\}$，则 $N \cup \complement_U M =$（    ）
\begin{choices}
  \item $\{2,3,5\}$
  \item $\{1,3,4\}$
  \item $\{1,2,4,5\}$
  \item $\{2,3,4,5\}$
\end{choices}
\begin{solution}
\textbf{答案：}A

\textbf{分析：}利用集合的交并补运算即可得解.

\textbf{详解：}因为全集 $U = \{1,2,3,4,5\}$，集合 $M = \{1,4\}$，所以 $\complement_U M = \{2,3,5\}$，又 $N = \{2,5\}$，所以 $N \cup \complement_U M = \{2,3,5\}$，故选A.
\end{solution}
\end{question}

\begin{question}[points=5]
$\frac{5(1+i^3)}{(2+i)(2-i)} =$（    ）
\begin{choices}
  \item $-1$
  \item $1$
  \item $1-i$
  \item $1+i$
\end{choices}
\begin{solution}
\textbf{答案：}C

\textbf{分析：}利用复数的四则运算求解即可.

\textbf{详解：}$\frac{5(1+i^3)}{(2+i)(2-i)} = \frac{5(1-i)}{5} = 1-i$，故选C.
\end{solution}
\end{question}

\begin{question}[points=5]
已知向量 $\vec{a} = (3,1), \vec{b} = (2,2)$，则 $\cos\langle \vec{a}+\vec{b}, \vec{a}-\vec{b} \rangle =$（    ）
\begin{choices}
  \item $\frac{1}{17}$
  \item $\frac{\sqrt{17}}{17}$
  \item $\frac{\sqrt{5}}{5}$
  \item $\frac{2\sqrt{5}}{5}$
\end{choices}
\begin{solution}
\textbf{答案：}B

\textbf{分析：}利用平面向量模与数量积的坐标表示分别求得 $|\vec{a}+\vec{b}|, |\vec{a}-\vec{b}|, (\vec{a}+\vec{b})\cdot(\vec{a}-\vec{b})$，从而利用平面向量余弦的运算公式即可得解.

\textbf{详解：}因为 $\vec{a} = (3,1), \vec{b} = (2,2)$，所以 $\vec{a}+\vec{b} = (5,3), \vec{a}-\vec{b} = (1,-1)$，则 $|\vec{a}+\vec{b}| = \sqrt{5^2+3^2} = \sqrt{34}, |\vec{a}-\vec{b}| = \sqrt{1+1} = \sqrt{2}, (\vec{a}+\vec{b})\cdot(\vec{a}-\vec{b}) = 5\times1+3\times(-1)=2$，所以 $\cos\langle \vec{a}+\vec{b}, \vec{a}-\vec{b} \rangle = \frac{2}{\sqrt{34}\times\sqrt{2}} = \frac{\sqrt{17}}{17}$，故选B.
\end{solution}
\end{question}

\begin{question}[points=5]
某校文艺部有4名学生，其中高一、高二年级各2名．从这4名学生中随机选2名组织校文艺汇演，则这2名学生来自不同年级的概率为（    ）
\begin{choices}
  \item $\frac{1}{6}$
  \item $\frac{1}{3}$
  \item $\frac{1}{2}$
  \item $\frac{2}{3}$
\end{choices}
\begin{solution}
\textbf{答案：}D

\textbf{分析：}利用古典概率的概率公式，结合组合的知识即可得解.

\textbf{详解：}依题意，从这4名学生中随机选2名组织校文艺汇演，总的基本事件有 $\mathrm{C}_4^2 = 6$ 件，其中这2名学生来自不同年级的基本事件有 $\mathrm{C}_2^1 \mathrm{C}_2^1 = 4$，所以这2名学生来自不同年级的概率为 $\frac{4}{6} = \frac{2}{3}$，故选D.
\end{solution}
\end{question}

\begin{question}[points=5]
记 $S_n$ 为等差数列 $\{a_n\}$ 的前 $n$ 项和．若 $a_2+a_6=10, a_4a_8=45$，则 $S_5 =$（    ）
\begin{choices}
  \item $25$
  \item $22$
  \item $20$
  \item $15$
\end{choices}
\begin{solution}
\textbf{答案：}C

\textbf{分析：}方法一：根据题意直接求出等差数列 $\{a_n\}$ 的公差和首项，再根据前n项和公式即可解出；方法二：根据等差数列的性质求出等差数列 $\{a_n\}$ 的公差，再根据前n项和公式的性质即可解出．

\textbf{详解：}方法一：设等差数列 $\{a_n\}$ 的公差为 $d$，首项为 $a_1$，依题意可得，$a_2+a_6 = a_1+d+a_1+5d=10$，即 $a_1+3d=5$，又 $a_4a_8 = (a_1+3d)(a_1+7d)=45$，解得：$d=1, a_1=2$，所以 $S_5 = 5a_1+\frac{5\times4}{2}\times d = 5\times2+10=20$．故选C.
方法二：$a_2+a_6=2a_4=10$，$a_4a_8=45$，所以 $a_4=5$，$a_8=9$，从而 $d=\frac{a_8-a_4}{8-4}=1$，于是 $a_3=a_4-d=5-1=4$，所以 $S_5=5a_3=20$．故选C.
\end{solution}
\end{question}

\begin{question}[points=5]
执行下边的程序框图，则输出的 $B =$（    ）

\begin{center}
  \includegraphics[width=.38\linewidth]{question-6-1.jpg}
\end{center}
\begin{choices}
  \item $21$
  \item $34$
  \item $55$
  \item $89$
\end{choices}
\begin{solution}
\textbf{答案：}B

\textbf{分析：}根据程序框图模拟运行即可解出．

\textbf{详解：}当 $k=1$ 时，判断框条件满足，第一次执行循环体，$A=1+2=3$，$B=3+2=5$，$k=1+1=2$；当 $k=2$ 时，判断框条件满足，第二次执行循环体，$A=3+5=8$，$B=8+5=13$，$k=2+1=3$；当 $k=3$ 时，判断框条件满足，第三次执行循环体，$A=8+13=21$，$B=21+13=34$，$k=3+1=4$；当 $k=4$ 时，判断框条件不满足，跳出循环体，输出 $B=34$．故选B.
\end{solution}
\end{question}

\begin{question}[points=5]
设 $F_1, F_2$ 为椭圆 $C: \frac{x^2}{5}+y^2=1$ 的两个焦点，点 $P$ 在 $C$ 上，若 $\overrightarrow{PF_1} \cdot \overrightarrow{PF_2}=0$，则 $PF_1 \cdot PF_2 =$（    ）
\begin{choices}
  \item $1$
  \item $2$
  \item $4$
  \item $5$
\end{choices}
\begin{solution}
\textbf{答案：}B

\textbf{分析：}方法一：根据焦点三角形面积公式求出 $\triangle PF_1F_2$ 的面积，即可解出；方法二：根据椭圆的定义以及勾股定理即可解出．

\textbf{详解：}方法一：因为 $\overrightarrow{PF_1} \cdot \overrightarrow{PF_2}=0$，所以 $\angle F_1PF_2=90^\circ$，从而 $S_{\triangle PF_1F_2} = b^2\tan45^\circ = 1 = \frac{1}{2} \cdot PF_1 \cdot PF_2$，所以 $PF_1 \cdot PF_2 = 2$．故选B.
方法二：因为 $\overrightarrow{PF_1} \cdot \overrightarrow{PF_2}=0$，所以 $\angle F_1PF_2=90^\circ$，由椭圆方程可知，$c^2=5-1=4 \Rightarrow c=2$，所以 $PF_1^2+PF_2^2 = F_1F_2^2 = 4^2=16$，又 $PF_1+PF_2=2a=2\sqrt{5}$，平方得：$PF_1^2+PF_2^2+2PF_1PF_2=16+2PF_1PF_2=20$，所以 $PF_1 \cdot PF_2=2$．故选B.
\end{solution}
\end{question}

\begin{question}[points=5]
曲线 $y = \frac{e^x}{x+1}$ 在点 $\left(1, \frac{e}{2}\right)$ 处的切线方程为（    ）
\begin{choices}
  \item $y = \frac{e}{4}x$
  \item $y = \frac{e}{2}x$
  \item $y = \frac{e}{4}x + \frac{e}{4}$
  \item $y = \frac{e}{2}x + \frac{3e}{4}$
\end{choices}
\begin{solution}
\textbf{答案：}C

\textbf{分析：}先由切点设切线方程，再求函数的导数，把切点的横坐标代入导数得到切线的斜率，代入所设方程即可求解.

\textbf{详解：}设曲线 $y = \frac{e^x}{x+1}$ 在点 $\left(1, \frac{e}{2}\right)$ 处的切线方程为 $y - \frac{e}{2} = k(x-1)$，因为 $y = \frac{e^x}{x+1}$，所以 $y' = \frac{e^x(x+1)-e^x}{(x+1)^2} = \frac{xe^x}{(x+1)^2}$，所以 $k = y'|_{x=1}= \frac{e}{4}$，所以 $y - \frac{e}{2} = \frac{e}{4}(x-1)$，即 $y = \frac{e}{4}x + \frac{e}{4}$，故选C.
\end{solution}
\end{question}

\begin{question}[points=5]
已知双曲线 $\frac{x^2}{a^2}-\frac{y^2}{b^2}=1(a>0,b>0)$ 的离心率为 $\sqrt{5}$，其中一条渐近线与圆 $(x-2)^2+(y-3)^2=1$ 交于A，B两点，则 $|AB| =$（    ）
\begin{choices}
  \item $\frac{\sqrt{5}}{5}$
  \item $\frac{2\sqrt{5}}{5}$
  \item $\frac{3\sqrt{5}}{5}$
  \item $\frac{4\sqrt{5}}{5}$
\end{choices}
\begin{solution}
\textbf{答案：}D

\textbf{分析：}根据离心率得出双曲线渐近线方程，再由圆心到直线的距离及圆半径可求弦长.

\textbf{详解：}由 $e=\sqrt{5}$，则 $\frac{c^2}{a^2}= \frac{a^2+b^2}{a^2}=1+\frac{b^2}{a^2}=5$，解得 $\frac{b}{a}=2$，所以双曲线的一条渐近线不妨取 $y=2x$，则圆心 $(2,3)$ 到渐近线的距离 $d = \frac{|2\times2-3|}{\sqrt{2^2+1}} = \frac{\sqrt{5}}{5}$，所以弦长 $|AB| = 2\sqrt{r^2-d^2} = 2\sqrt{1-\frac{1}{5}} = \frac{4\sqrt{5}}{5}$，故选D.
\end{solution}
\end{question}

\begin{question}[points=5]
在三棱锥 $P-ABC$ 中，$\triangle ABC$ 是边长为2的等边三角形，$PA=PB=2, PC=\sqrt{6}$，则该棱锥的体积为（    ）
\begin{choices}
  \item $1$
  \item $\sqrt{3}$
  \item $2$
  \item $3$
\end{choices}
\begin{solution}
\textbf{答案：}A

\textbf{分析：}证明 $AB \perp$ 平面 $PEC$，分割三棱锥为共底面两个小三棱锥，其高之和为AB得解.

\textbf{详解：}取AB中点E，连接PE,CE，如图，$\because \triangle ABC$是边长为2的等边三角形，$PA=PB=2$，$\therefore PE \perp AB, CE \perp AB$，又$PE, CE \subset$平面$PEC$，$PE \cap CE = E$，$\therefore AB \perp$平面$PEC$，又$PE=CE=2\times\frac{\sqrt{3}}{2}=\sqrt{3}$，$PC=\sqrt{6}$，故$PC^2=PE^2+CE^2$，即$PE \perp CE$，所以$V=V_{B-PEC}+V_{A-PEC}=\frac{1}{3}S_{\triangle PEC}\cdot AB = \frac{1}{3}\times\frac{1}{2}\times\sqrt{3}\times\sqrt{3}\times2=1$，故选A.
\end{solution}
\end{question}

\begin{question}[points=5]
已知函数 $f(x)=e^x-(x-1)^2$．记 $a=f\left(\frac{\sqrt{2}}{2}\right), b=f\left(\frac{\sqrt{3}}{2}\right), c=f\left(\frac{\sqrt{6}}{2}\right)$，则（    ）
\begin{choices}
  \item $b>c>a$
  \item $b>a>c$
  \item $c>b>a$
  \item $c>a>b$
\end{choices}
\begin{solution}
\textbf{答案：}A

\textbf{分析：}利用作差法比较自变量的大小，再根据指数函数的单调性及二次函数的性质判断即可.

\textbf{详解：}令 $g(x)=-(x-1)^2$，则$g(x)$开口向下，对称轴为$x=1$，因为$\frac{\sqrt{6}}{2}-1-(1-\frac{\sqrt{3}}{2})=\frac{\sqrt{6}+\sqrt{3}}{2}-\frac{4}{2}$，而$(\sqrt{6}+\sqrt{3})^2-4^2=9+6\sqrt{2}-16=6\sqrt{2}-7>0$，所以$\frac{\sqrt{6}}{2}-1>1-\frac{\sqrt{3}}{2}$，由二次函数性质知$g(\frac{\sqrt{6}}{2})<g(\frac{\sqrt{3}}{2})$；因为$\frac{\sqrt{6}}{2}-1-(1-\frac{\sqrt{2}}{2})=\frac{\sqrt{6}+\sqrt{2}}{2}-\frac{4}{2}$，而$(\sqrt{6}+\sqrt{2})^2-4^2=8+4\sqrt{3}-16=4\sqrt{3}-8<0$，即$\frac{\sqrt{6}}{2}-1<1-\frac{\sqrt{2}}{2}$，所以$g(\frac{\sqrt{6}}{2})>g(\frac{\sqrt{2}}{2})$，综上，$g(\frac{\sqrt{2}}{2})<g(\frac{\sqrt{6}}{2})<g(\frac{\sqrt{3}}{2})$，又$y=e^x$为增函数，故$a<c<b$，即$b>c>a$，故选A.
\end{solution}
\end{question}

\begin{question}[points=5]
函数 $y=f(x)$ 的图象由 $y=\cos\left(2x+\frac{\pi}{6}\right)$ 的图象向左平移 $\frac{\pi}{6}$ 个单位长度得到，则 $y=f(x)$ 的图象与直线 $y=\frac{1}{2}x-\frac{1}{2}$ 的交点个数为（    ）
\begin{choices}
  \item $1$
  \item $2$
  \item $3$
  \item $4$
\end{choices}
\begin{solution}
\textbf{答案：}C

\textbf{分析：}先利用三角函数平移的性质求得 $f(x)=-\sin2x$，再作出 $f(x)$ 与 $y=\frac{1}{2}x-\frac{1}{2}$ 的部分大致图像，考虑特殊点处 $f(x)$ 与 $y=\frac{1}{2}x-\frac{1}{2}$ 的大小关系，从而精确图像，由此得解.

\textbf{详解：}因为 $y=\cos\left(2x+\frac{\pi}{6}\right)$ 向左平移 $\frac{\pi}{6}$ 个单位所得函数为 $y=\cos\left[2(x+\frac{\pi}{6})+\frac{\pi}{6}\right]=\cos\left(2x+\frac{\pi}{2}\right)=-\sin2x$，所以 $f(x)=-\sin2x$，而 $y=\frac{1}{2}x-\frac{1}{2}$ 显然过 $(0,-\frac{1}{2})$ 与 $(1,0)$ 两点，作出 $f(x)$ 与 $y=\frac{1}{2}x-\frac{1}{2}$ 的部分大致图像，考虑 $2x=-\frac{3\pi}{2},2x=\frac{3\pi}{2},2x=\frac{7\pi}{2}$，即 $x=-\frac{3\pi}{4},x=\frac{3\pi}{4},x=\frac{7\pi}{4}$ 处 $f(x)$ 与 $y=\frac{1}{2}x-\frac{1}{2}$ 的大小关系，当 $x=-\frac{3\pi}{4}$ 时，$f(-\frac{3\pi}{4})=-\sin(-\frac{3\pi}{2})=-1$，$y=\frac{1}{2}\times(-\frac{3\pi}{4})-\frac{1}{2}=-\frac{3\pi+4}{8}<-1$；当 $x=\frac{3\pi}{4}$ 时，$f(\frac{3\pi}{4})=-\sin\frac{3\pi}{2}=1$，$y=\frac{1}{2}\times\frac{3\pi}{4}-\frac{1}{2}=\frac{3\pi-4}{8}<1$；当 $x=\frac{7\pi}{4}$ 时，$f(\frac{7\pi}{4})=-\sin\frac{7\pi}{2}=1$，$y=\frac{1}{2}\times\frac{7\pi}{4}-\frac{1}{2}=\frac{7\pi-4}{8}>1$；所以由图可知，$f(x)$ 与 $y=\frac{1}{2}x-\frac{1}{2}$ 的交点个数为3，故选C.
\end{solution}
\end{question}

\section*{二、填空题}
本大题共4小题，每小题5分，共20分.

\begin{question}[points=5]
记 $S_n$ 为等比数列 $\{a_n\}$ 的前 $n$ 项和．若 $8S_6=7S_3$，则 $\{a_n\}$ 的公比为．
\fillin[{$-\frac{1}{2}$}]
\begin{solution}
\textbf{答案：}$-\frac{1}{2}$

\textbf{分析：}先分析 $q\neq1$，再由等比数列的前n项和公式和平方差公式化简即可求出公比 $q$.

\textbf{详解：}若 $q=1$，则由 $8S_6=7S_3$ 得 $8\cdot6a_1=7\cdot3a_1$，则 $a_1=0$，不合题意，所以 $q\neq1$．当 $q\neq1$ 时，因为 $8S_6=7S_3$，所以 $8\cdot\frac{a_1(1-q^6)}{1-q}=7\cdot\frac{a_1(1-q^3)}{1-q}$，即 $8(1-q^6)=7(1-q^3)$，即 $8(1+q^3)(1-q^3)=7(1-q^3)$，即 $8(1+q^3)=7$，解得 $q^3=-\frac{1}{8}$，所以 $q=-\frac{1}{2}$．
\end{solution}
\end{question}

\begin{question}[points=5]
若 $f(x)=(x-1)^2+ax+\sin\left(x+\frac{\pi}{2}\right)$ 为偶函数，则 $a=$ ．
\fillin[{$2$}]
\begin{solution}
\textbf{答案：}$2$

\textbf{分析：}根据常见函数的奇偶性直接求解即可.

\textbf{详解：}因为 $f(x)=(x-1)^2+ax+\sin\left(x+\frac{\pi}{2}\right) = (x-1)^2+ax+\cos x = x^2+(a-2)x+1+\cos x$，且函数为偶函数，所以 $a-2=0$，解得 $a=2$．
\end{solution}
\end{question}

\begin{question}[points=5]
若 $x,y$ 满足约束条件 $\begin{cases}3x-2y\le3\\-2x+3y\le3\\x+y\ge1\end{cases}$，则 $z=3x+2y$ 的最大值为．
\fillin[{$15$}]
\begin{solution}
\textbf{答案：}$15$

\textbf{分析：}由约束条件作出可行域，根据线性规划求最值即可.

\textbf{详解：}作出可行域，如图，由图可知，当目标函数 $y=-\frac{3}{2}x+\frac{z}{2}$ 过点 $A$ 时，$z$ 有最大值，由 $\begin{cases}-2x+3y=3\\3x-2y=3\end{cases}$ 可得 $\begin{cases}x=3\\y=3\end{cases}$，即 $A(3,3)$，所以 $z_{\max}=3\times3+2\times3=15$．
\end{solution}
\end{question}

\begin{question}[points=5]
在正方体 $ABCD-A_1B_1C_1D_1$ 中，$AB=4$，$O$ 为 $AC_1$ 的中点，若该正方体的棱与球 $O$ 的球面有公共点，则球 $O$ 的半径的取值范围是．
\fillin[{$[2\sqrt{2}, 2\sqrt{3}]$}]
\begin{solution}
\textbf{答案：}$[2\sqrt{2}, 2\sqrt{3}]$

\textbf{分析：}当球是正方体的外接球时半径最大，当边长为4的正方形是球的大圆的内接正方形时半径达到最小.

\textbf{详解：}设球的半径为 $R$．当球是正方体的外接球时，恰好经过正方体的每个顶点，所求的球的半径最大，若半径变得更大，球会包含正方体，导致球面和棱没有交点，正方体的外接球直径 $2R'$ 为体对角线长 $AC_1=\sqrt{4^2+4^2+4^2}=4\sqrt{3}$，即 $2R'=4\sqrt{3}, R'=2\sqrt{3}$，故 $R_{\max}=2\sqrt{3}$；分别取侧棱 $AA_1, BB_1, CC_1, DD_1$ 的中点 $M,H,G,N$，显然四边形 $MNGH$ 是边长为4的正方形，且 $O$ 为正方形 $MNGH$ 的对角线交点，连接 $MG$，则 $MG=4\sqrt{2}$，当球的一个大圆恰好是四边形 $MNGH$ 的外接圆，球的半径达到最小，即 $R$ 的最小值为 $2\sqrt{2}$．综上，$R\in[2\sqrt{2},2\sqrt{3}]$．
\end{solution}
\end{question}

\section*{三、解答题}
共70分.解答应写出文字说明、证明过程或演算步骤.第17~21题为必考题，每个试题考生都必须作答.第22、23题为选考题，考生根据要求作答.（一）必考题：共60分.

\begin{question}
记 $\triangle ABC$ 的内角 $A,B,C$ 的对边分别为 $a,b,c$，已知 $\frac{b^2+c^2-a^2}{\cos A}=2$．

（1）求 $bc$；

（2）若 $\frac{a\cos B-b\cos A}{a\cos B+b\cos A}-\frac{b}{c}=1$，求 $\triangle ABC$ 面积．
\begin{solution}
\textbf{答案：}（1）$bc=1$；（2）$\frac{\sqrt{3}}{4}$

\textbf{分析：}（1）根据余弦定理即可解出；（2）由（1）可知，只需求出 $\sin A$ 即可得到三角形面积，对等式恒等变换，即可解出．

\textbf{详解：}（1）因为 $a^2=b^2+c^2-2bc\cos A$，所以 $\frac{b^2+c^2-a^2}{\cos A}=\frac{2bc\cos A}{\cos A}=2bc=2$，解得 $bc=1$．
（2）由正弦定理可得 $\frac{a\cos B-b\cos A}{a\cos B+b\cos A}-\frac{b}{c}=\frac{\sin A\cos B-\sin B\cos A}{\sin A\cos B+\sin B\cos A}-\frac{\sin B}{\sin C}=\frac{\sin(A-B)}{\sin(A+B)}-\frac{\sin B}{\sin(A+B)}=\frac{\sin(A-B)-\sin B}{\sin(A+B)}=1$，变形可得 $\sin(A-B)-\sin(A+B)=\sin B$，即 $-2\cos A\sin B=\sin B$，而 $0<\sin B\le1$，所以 $\cos A=-\frac{1}{2}$，又 $0<A<\pi$，所以 $\sin A=\frac{\sqrt{3}}{2}$，故 $\triangle ABC$ 的面积为 $S_{\triangle ABC}=\frac{1}{2}bc\sin A=\frac{1}{2}\times1\times\frac{\sqrt{3}}{2}=\frac{\sqrt{3}}{4}$．
\end{solution}
\end{question}

\begin{question}
如图，在三棱柱 $ABC-A_1B_1C_1$ 中，$A_1C\perp$ 平面 $ABC$，$\angle ACB=90^\circ$．

（1）证明：平面 $ACC_1A_1\perp$ 平面 $BB_1C_1C$；

（2）设 $AB=A_1B$，$AA_1=2$，求四棱锥 $A_1-BB_1C_1C$ 的高．
\begin{solution}
\textbf{答案：}（1）证明见解析；（2）高为1

\textbf{分析：}（1）由 $A_1C\perp$ 平面 $ABC$ 得 $A_1C\perp BC$，又因为 $AC\perp BC$，可证 $BC\perp$ 平面 $ACC_1A_1$，从而证得平面 $ACC_1A_1\perp$ 平面 $BCC_1B_1$；（2）过点 $A_1$ 作 $A_1O\perp CC_1$，可证四棱锥的高为 $A_1O$，由三角形全等可证 $A_1C=AC$，从而证得 $O$ 为 $CC_1$ 中点，设 $A_1C=AC=x$，由勾股定理可求出 $x$，再由勾股定理即可求 $A_1O$．

\textbf{详解：}（1）证明：因为 $A_1C\perp$ 平面 $ABC$，$BC\subset$ 平面 $ABC$，所以 $A_1C\perp BC$，又因为 $\angle ACB=90^\circ$，即 $AC\perp BC$，$A_1C, AC\subset$ 平面 $ACC_1A_1$，$A_1C\cap AC=C$，所以 $BC\perp$ 平面 $ACC_1A_1$，又因为 $BC\subset$ 平面 $BCC_1B_1$，所以平面 $ACC_1A_1\perp$ 平面 $BCC_1B_1$．
（2）如图，过点 $A_1$ 作 $A_1O\perp CC_1$，垂足为 $O$．因为平面 $ACC_1A_1\perp$ 平面 $BCC_1B_1$，平面 $ACC_1A_1\cap$ 平面 $BCC_1B_1=CC_1$，$A_1O\subset$ 平面 $ACC_1A_1$，所以 $A_1O\perp$ 平面 $BCC_1B_1$，所以四棱锥 $A_1-BB_1C_1C$ 的高为 $A_1O$．因为 $A_1C\perp$ 平面 $ABC$，$AC, BC\subset$ 平面 $ABC$，所以 $A_1C\perp BC, A_1C\perp AC$，又因为 $A_1B=AB$，$BC$ 为公共边，所以 $\triangle ABC$ 与 $\triangle A_1BC$ 全等，所以 $A_1C=AC$．设 $A_1C=AC=x$，则 $A_1C_1=x$，所以 $O$ 为 $CC_1$ 中点，$OC_1=\frac{1}{2}AA_1=1$，又因为 $A_1C\perp AC$，所以 $A_1C^2+AC^2=AA_1^2$，即 $x^2+x^2=2^2$，解得 $x=\sqrt{2}$，所以 $A_1O=\sqrt{A_1C_1^2-OC_1^2}=\sqrt{(\sqrt{2})^2-1^2}=1$．所以四棱锥 $A_1-BB_1C_1C$ 的高为1．
\end{solution}
\end{question}

\begin{question}
一项试验旨在研究臭氧效应，试验方案如下：选40只小白鼠，随机地将其中20只分配到试验组，另外20只分配到对照组，试验组的小白鼠饲养在高浓度臭氧环境，对照组的小白鼠饲养在正常环境，一段时间后统计每只小白鼠体重的增加量（单位：g）．试验结果如下：

对照组的小白鼠体重的增加量从小到大排序为

15.2 18.8 20.2 21.3 22.5 23.2 25.8 26.5 27.5 30.1

32.6 34.3 34.8 35.6 35.6 35.8 36.2 37.3 40.5 43.2

试验组的小白鼠体重的增加量从小到大排序为

7.8 9.2 11.4 12.4 13.2 15.5 16.5 18.0 18.8 19.2

19.8 20.2 21.6 22.8 23.6 23.9 25.1 28.2 32.3 36.5

（1）计算试验组的样本平均数；

（2）（ⅰ）求40只小白鼠体重的增加量的中位数 $m$，再分别统计两样本中小于 $m$ 与不小于 $m$ 的数据的个数，完成如下列联表

$\begin{array}{c|cc} & <m & \ge m \\ \hline \text{对照组} & & \\ \text{试验组} & & \\ \end{array}$

（ⅱ）根据（i）中的列联表，能否有95%的把握认为小白鼠在高浓度臭氧环境中与在正常环境中体重的增加量有差异？

附：$K^2=\frac{n(ad-bc)^2}{(a+b)(c+d)(a+c)(b+d)}$，$P(K^2\ge k)$ 0.100 0.050 0.010，$k$ 2.706 3.841 6.635
\begin{solution}
\textbf{答案：}（1）19.8；（2）（i）$m=23.4$，列联表见解析；（ii）能有95%的把握认为有差异

\textbf{分析：}（1）直接根据均值定义求解；（2）（i）根据中位数的定义即可求得 $m=23.4$，从而求得列联表；（ii）利用独立性检验的卡方计算进行检验，即可得解.

\textbf{详解：}（1）试验组样本平均数为：$\frac{1}{20}(7.8+9.2+11.4+12.4+13.2+15.5+16.5+18.0+18.8+19.2+19.8+20.2+21.6+22.8+23.6+23.9+25.1+28.2+32.3+36.5)=\frac{396}{20}=19.8$．
（2）（i）依题意，可知这40只小鼠体重的中位数是将两组数据合在一起，从小到大排后第20位与第21位数据的平均数，由原数据可得第11位数据为18.8，后续依次为19.2,19.8,20.2,20.2,21.3,21.6,22.5,22.8,23.2,23.6,$\cdots$，故第20位为23.2，第21位数据为23.6，所以 $m=\frac{23.2+23.6}{2}=23.4$，故列联表为：
$\begin{array}{c|cc|c} & <m & \ge m & \text{合计} \\ \hline \text{对照组} & 6 & 14 & 20 \\ \text{试验组} & 14 & 6 & 20 \\ \hline \text{合计} & 20 & 20 & 40 \end{array}$
（ii）由（i）可得，$K^2=\frac{40\times(6\times6-14\times14)^2}{20\times20\times20\times20}=6.400>3.841$，所以能有95%的把握认为小白鼠在高浓度臭氧环境中与在正常环境中体重的增加量有差异.
\end{solution}
\end{question}

\begin{question}
已知函数 $f(x)=ax-\frac{\sin x}{\cos^2 x}, x\in(0,\frac{\pi}{2})$．

（1）当 $a=1$ 时，讨论 $f(x)$ 的单调性；

（2）若 $f(x)+\sin x<0$，求 $a$ 的取值范围．
\begin{solution}
\textbf{答案：}（1）$f(x)$ 在 $(0,\frac{\pi}{2})$ 上单调递减；（2）$a\le0$

\textbf{分析：}（1）代入 $a=1$ 后，再对 $f(x)$ 求导，同时利用三角函数的平方关系化简 $f'(x)$，再利用换元法判断得其分子与分母的正负情况，从而得解；（2）法一：构造函数 $g(x)=f(x)+\sin x$，从而得到 $g(x)<0$，注意到 $g(0)=0$，从而得到 $g'(0)\le0$，进而得到 $a\le0$，再分类讨论 $a=0$ 与 $a<0$ 两种情况即可得解；法二：先化简并判断得 $\sin x-\frac{\sin x}{\cos^2 x}<0$ 恒成立，再分类讨论 $a=0$，$a<0$ 与 $a>0$ 三种情况，利用零点存在定理与隐零点的知识判断得 $a>0$ 时不满足题意，从而得解.

\textbf{详解：}（1）因为 $a=1$，所以 $f(x)=x-\frac{\sin x}{\cos^2 x}, x\in(0,\frac{\pi}{2})$，则 $f'(x)=1-\frac{\cos x\cdot\cos^2 x-\cos x\cdot(-\sin x)\sin x}{\cos^4 x}=1-\frac{\cos^2 x+2\sin^2 x}{\cos^3 x}=\frac{\cos^3 x-\cos^2 x-2(1-\cos^2 x)}{\cos^3 x}=\frac{\cos^3 x+\cos^2 x-2}{\cos^3 x}$，令 $t=\cos x$，由于 $x\in(0,\frac{\pi}{2})$，所以 $t=\cos x\in(0,1)$，所以 $\cos^3 x+\cos^2 x-2=t^3+t^2-2=(t-1)(t^2+2t+2)$，因为 $t^2+2t+2>0$，$t-1<0$，$\cos^3 x=t^3>0$，所以 $f'(x)<0$ 在 $(0,\frac{\pi}{2})$ 上恒成立，所以 $f(x)$ 在 $(0,\frac{\pi}{2})$ 上单调递减.
（2）法一：构建 $g(x)=f(x)+\sin x=ax-\frac{\sin x}{\cos^2 x}+\sin x(0<x<\frac{\pi}{2})$，则 $g'(x)=a-\frac{1+\sin^2 x}{\cos^3 x}+\cos x$，若 $g(x)=f(x)+\sin x<0$，且 $g(0)=f(0)+\sin0=0$，则 $g'(0)=a-1+1=a\le0$，解得 $a\le0$，当 $a=0$ 时，因为 $\sin x-\frac{\sin x}{\cos^2 x}=\sin x\left(1-\frac{1}{\cos^2 x}\right)$，又 $x\in(0,\frac{\pi}{2})$，所以 $0<\sin x<1$，$0<\cos x<1$，则 $\frac{1}{\cos^2 x}>1$，所以 $f(x)+\sin x=\sin x-\frac{\sin x}{\cos^2 x}<0$，满足题意；当 $a<0$ 时，由于 $0<x<\frac{\pi}{2}$，显然 $ax<0$，所以 $f(x)+\sin x=ax-\frac{\sin x}{\cos^2 x}+\sin x<\sin x-\frac{\sin x}{\cos^2 x}<0$，满足题意；综上所述：若 $f(x)+\sin x<0$，等价于 $a\le0$，所以 $a$ 的取值范围为 $(-\infty,0]$．
法二：因为 $\sin x-\frac{\sin x}{\cos^2 x}=\frac{\sin x\cos^2 x-\sin x}{\cos^2 x}=\frac{-\sin^3 x}{\cos^2 x}$，因为 $x\in(0,\frac{\pi}{2})$，所以 $0<\sin x<1$，$0<\cos x<1$，故 $\sin x-\frac{\sin x}{\cos^2 x}<0$ 在 $(0,\frac{\pi}{2})$ 上恒成立，所以当 $a=0$ 时，$f(x)+\sin x=\sin x-\frac{\sin x}{\cos^2 x}<0$，满足题意；当 $a<0$ 时，由于 $0<x<\frac{\pi}{2}$，显然 $ax<0$，所以 $f(x)+\sin x=ax-\frac{\sin x}{\cos^2 x}+\sin x<\sin x-\frac{\sin x}{\cos^2 x}<0$，满足题意；当 $a>0$ 时，因为 $f(x)+\sin x=ax-\frac{\sin x}{\cos^2 x}+\sin x=ax-\frac{\sin^3 x}{\cos^2 x}$，令 $g(x)=ax-\frac{\sin^3 x}{\cos^2 x}(0<x<\frac{\pi}{2})$，则 $g'(x)=a-\frac{3\sin^2 x\cos^2 x+2\sin^4 x}{\cos^3 x}$，注意到 $g'(0)=a-0=a>0$，若 $\forall0<x<\frac{\pi}{2}$，$g'(x)>0$，则 $g(x)$ 在 $(0,\frac{\pi}{2})$ 上单调递增，注意到 $g(0)=0$，所以 $g(x)>g(0)=0$，即 $f(x)+\sin x>0$，不满足题意；若 $\exists0<x_0<\frac{\pi}{2}$，$g'(x_0)<0$，则 $g'(0)g'(x_0)<0$，所以在 $(0,\frac{\pi}{2})$ 上最靠近 $x=0$ 处必存在零点 $x_1\in(0,\frac{\pi}{2})$，使得 $g'(x_1)=0$，此时 $g'(x)$ 在 $(0,x_1)$ 上有 $g'(x)>0$，所以 $g(x)$ 在 $(0,x_1)$ 上单调递增，则在 $(0,x_1)$ 上有 $g(x)>g(0)=0$，即 $f(x)+\sin x>0$，不满足题意；综上：$a\le0$.
\end{solution}
\end{question}

\begin{question}
已知直线 $x-2y+1=0$ 与抛物线 $C: y^2=2px(p>0)$ 交于 $A,B$ 两点，$|AB|=4\sqrt{15}$．

（1）求 $p$；

（2）设 $F$ 为 $C$ 的焦点，$M,N$ 为 $C$ 上两点，且 $\overrightarrow{FM}\cdot\overrightarrow{FN}=0$，求 $\triangle MFN$ 面积的最小值．
\begin{solution}
\textbf{答案：}（1）$p=2$；（2）$12-8\sqrt{2}$

\textbf{分析：}（1）利用直线与抛物线的位置关系，联立直线和抛物线方程求出弦长即可得出 $p$；（2）设直线 $MN: x=my+n$，$M(x_1,y_1),N(x_2,y_2)$，利用 $\overrightarrow{MF}\cdot\overrightarrow{NF}=0$，找到 $m,n$ 的关系，以及 $\triangle MNF$ 的面积表达式，再结合函数的性质即可求出其最小值．

\textbf{详解：}（1）设 $A(x_A,y_A), B(x_B,y_B)$，由 $\begin{cases}x-2y+1=0\\y^2=2px\end{cases}$ 可得，$y^2-4py+2p=0$，所以 $y_A+y_B=4p, y_Ay_B=2p$，所以 $|AB|=\sqrt{(x_A-x_B)^2+(y_A-y_B)^2}=\sqrt{5}|y_A-y_B|=\sqrt{5}\times\sqrt{(y_A+y_B)^2-4y_Ay_B}=4\sqrt{15}$，即 $2p^2-p-6=0$，因为 $p>0$，解得：$p=2$．
（2）因为 $F(1,0)$，显然直线 $MN$ 的斜率不可能为零，设直线 $MN: x=my+n$，$M(x_1,y_1), N(x_2,y_2)$，由 $\begin{cases}y^2=4x\\x=my+n\end{cases}$ 可得，$y^2-4my-4n=0$，所以 $y_1+y_2=4m, y_1y_2=-4n$，$\Delta=16m^2+16n>0\Rightarrow m^2+n>0$，因为 $\overrightarrow{MF}\cdot\overrightarrow{NF}=0$，所以 $(x_1-1)(x_2-1)+y_1y_2=0$，即 $(my_1+n-1)(my_2+n-1)+y_1y_2=0$，亦即 $(m^2+1)y_1y_2+m(n-1)(y_1+y_2)+(n-1)^2=0$，将 $y_1+y_2=4m, y_1y_2=-4n$ 代入得，$4m^2=n^2-6n+1$，$4(m^2+n)=(n-1)^2>0$，所以 $n\neq1$，且 $n^2-6n+1\ge0$，解得 $n\ge3+2\sqrt{2}$ 或 $n\le3-2\sqrt{2}$．设点 $F$ 到直线 $MN$ 的距离为 $d$，所以 $d=\frac{|n-1|}{\sqrt{1+m^2}}$，$|MN|=\sqrt{(x_1-x_2)^2+(y_1-y_2)^2}=\sqrt{1+m^2}|y_1-y_2|=\sqrt{1+m^2}\sqrt{16m^2+16n}=2\sqrt{1+m^2}\sqrt{4(n^2-6n+1)+16n}=2\sqrt{1+m^2}|n-1|$，所以 $\triangle MNF$ 的面积 $S=\frac{1}{2}\times|MN|\times d=\frac{1}{2}\times2\sqrt{1+m^2}|n-1|\times\frac{|n-1|}{\sqrt{1+m^2}}=(n-1)^2$，而 $n\ge3+2\sqrt{2}$ 或 $n\le3-2\sqrt{2}$，所以当 $n=3-2\sqrt{2}$ 时，$\triangle MNF$ 的面积 $S_{\min}=(2-2\sqrt{2})^2=12-8\sqrt{2}$．
\end{solution}
\end{question}

\section*{（二）选考题}
共10分.请考生在第22、23题中任选一题作答.如果多做，则按所做的第一题计分.

\begin{question}[points=10]
[选修4-4：坐标系与参数方程]（10分）

已知点 $P(2,1)$，直线 $l: \begin{cases}x=2+t\cos\alpha\\ y=1+t\sin\alpha\end{cases}$（$t$ 为参数），$\alpha$ 为 $l$ 的倾斜角，$l$ 与 $x$ 轴正半轴、$y$ 轴正半轴分别交于 $A,B$，且 $PA\cdot PB=4$．

（1）求 $\alpha$；

（2）以坐标原点为极点，$x$ 轴正半轴为极轴建立极坐标系，求 $l$ 的极坐标方程．
\begin{solution}
\textbf{答案：}（1）$\frac{3\pi}{4}$；（2）$\rho\cos\theta+\rho\sin\theta-3=0$

\textbf{分析：}（1）根据 $t$ 的几何意义即可解出；（2）求出直线 $l$ 的普通方程，再根据直角坐标和极坐标互化公式即可解出．

\textbf{详解：}（1）因为 $l$ 与 $x$ 轴，$y$ 轴正半轴交于 $A,B$ 两点，所以 $\frac{\pi}{2}<\alpha<\pi$，令 $x=0$，$t_1=-\frac{2}{\cos\alpha}$，令 $y=0$，$t_2=-\frac{1}{\sin\alpha}$，所以 $PA\cdot PB=|t_2t_1|=\frac{2}{|\sin\alpha\cos\alpha|}=\frac{4}{|\sin2\alpha|}=4$，所以 $\sin2\alpha=\pm1$，即 $2\alpha=\frac{\pi}{2}+k\pi$，解得 $\alpha=\frac{\pi}{4}+\frac{k\pi}{2},k\in\mathbb{Z}$，因为 $\frac{\pi}{2}<\alpha<\pi$，所以 $\alpha=\frac{3\pi}{4}$．
（2）由（1）可知，直线 $l$ 的斜率为 $\tan\alpha=-1$，且过点 $(2,1)$，所以直线 $l$ 的普通方程为 $y-1=-(x-2)$，即 $x+y-3=0$，由 $x=\rho\cos\theta, y=\rho\sin\theta$ 可得直线 $l$ 的极坐标方程为 $\rho\cos\theta+\rho\sin\theta-3=0$．
\end{solution}
\end{question}

\begin{question}[points=10]
[选修4-5：不等式选讲]（10分）

已知 $f(x)=2|x-a|-a, a>0$．

（1）求不等式 $f(x)<x$ 的解集；

（2）若曲线 $y=f(x)$ 与 $x$ 轴所围成的图形的面积为2，求 $a$．
\begin{solution}
\textbf{答案：}（1）$(\frac{a}{3},3a)$；（2）$\frac{2\sqrt{6}}{3}$

\textbf{分析：}（1）分 $x\le a$ 和 $x>a$ 讨论即可；（2）写出分段函数，画出草图，表达面积解方程即可.

\textbf{详解：}（1）若 $x\le a$，则 $f(x)=2a-2x-a<x$，即 $3x>a$，解得 $x>\frac{a}{3}$，即 $\frac{a}{3}<x\le a$；若 $x>a$，则 $f(x)=2x-2a-a<x$，解得 $x<3a$，即 $a<x<3a$；综上，不等式的解集为 $(\frac{a}{3},3a)$．
（2）$f(x)=\begin{cases}-2x+a, & x\le a\\ 2x-3a, & x>a\end{cases}$，画出 $f(x)$ 的草图，则 $f(x)$ 与坐标轴围成 $\triangle ADO$ 与 $\triangle ABC$，$\triangle ABC$ 的高为 $a$，$D(0,a)$，$A(\frac{a}{2},0)$，$B(\frac{3a}{2},0)$，所以 $|AB|=a$，所以 $S_{\triangle OAD}+S_{\triangle ABC}=\frac{1}{2}\cdot OA\cdot a+\frac{1}{2}\cdot AB\cdot a=\frac{3}{4}a^2=2$，解得 $a=\frac{2\sqrt{6}}{3}$．
\end{solution}
\end{question}

